\section{Resumes of Key Individuals}

\subsection{David Daeyoung Lee, Principal Investigator}

\noindent\textit{Note: Publications and records prior to 2024 may appear under the name Dae Young Lee.}

\subsubsection{Education}

\begin{itemize}
  \item Ph.D., Aerospace Engineering, University of Michigan, 2016.
  \item M.S., Aerospace Engineering, University of Michigan, 2015.
  \item M.S., Mechanical Engineering, Pusan National University, South Korea, 2001.
  \item B.S., Mechanical Engineering, Pusan National University, South Korea, 1997.
\end{itemize}

\subsubsection{Professional Appointments}

\begin{itemize}
  \item Assistant Professor, California State University, Long Beach, Long Beach, CA, August 2026--present.
  \item Assistant Professor, Iowa State University, Ames, IA, August 2018--July 2026.
  \item Postdoctoral Researcher, The University of Texas at Austin, Austin, TX, July 2016--July 2018.
  \item Visiting Scholar, Naval Postgraduate School, Monterey, CA, May 2016--June 2016.
  \item Research Assistant, University of Michigan, Ann Arbor, MI, May 2010--April 2016.
  \item Project Manager, LG Industrial Systems, Seoul, Korea, May 2002--June 2009.
  \item Control Engineer, Hyundai Heavy Industry, Yong-in, Korea, November 2000--May 2002.
\end{itemize}

\subsubsection{Research Expertise}

Dr. Lee's research background is directly relevant to the proposed effort in hypersonic guidance, predictive control, trajectory simulation, and reinforcement learning-augmented guidance development. His expertise includes spacecraft dynamics and control, nonlinear model predictive control, optimal control, Mars entry, descent, and landing, CubeSat systems, and aerospace system optimization. This background provides a strong technical foundation for the proposed multi-model numerical predictor-corrector guidance framework and its extension to adaptive reinforcement learning-based guidance policy development.

\subsubsection{Selected Journal Publications}

\begin{enumerate}
  \item Y. Lee, D. D. Lee, and B. Wie, ``Autonomous numerical predictor-corrector guidance for human Mars landing missions,'' \textit{Aerospace Science and Technology}, vol. 156, Art. no. 109755, 2025.
  \item Y. Lee, D. Y. Lee, and B. Wie, ``A comparative study of entry guidance for Mars robotic and human landing missions,'' \textit{Acta Astronautica}, vol. 215, pp. 178--193, 2024.
  \item Y. Lee and D. Y. Lee, ``Entry trajectory generation for Mars robotic and human missions based on a predetermined bank angle profile,'' \textit{Advances in Space Research}, vol. 71, no. 8, pp. 3301--3312, 2023.
  \item Y. Lee, D. Y. Lee, S. H. Lee, and Y. Kim, ``A comparative study on model predictive control design for highway car-following scenarios: Space-domain and time-domain model,'' \textit{IEEE Access}, vol. 9, pp. 162291--162305, 2021.
  \item Y. He, Y. Kim, D. Y. Lee, and S.-H. Kim, ``Defensive ecological adaptive cruise control considering neighboring vehicles' blind-spot zones,'' \textit{IEEE Access}, vol. 9, pp. 152275--152287, 2021.
  \item J. Y. Kim, D. Y. Lee, J. Lee, and S. H. Lee, ``Parameter optimization of hybrid-tandem gas metal arc welding using analysis of variance-based Gaussian process regression,'' \textit{Metals}, vol. 11, no. 7, Art. no. 1087, 2021.
  \item G. Cervettini, S. Pastorelli, H. Park, D. Y. Lee, and M. Romano, ``Development and experimentation of a CubeSat magnetic attitude control system testbed,'' \textit{IEEE Transactions on Aerospace and Electronic Systems}, vol. 57, no. 2, pp. 1345--1350, 2020.
  \item D. Y. Lee, L. Leifsson, J.-Y. Kim, and S. H. Lee, ``Optimisation of hybrid tandem metal active gas welding using Gaussian process regression,'' \textit{Science and Technology of Welding and Joining}, vol. 25, no. 3, pp. 208--217, 2020.
  \item M. Mammarella, D. Y. Lee, H. Park, E. Capello, M. Dentis, and G. Guglieri, ``Attitude control of a small spacecraft via tube-based model predictive control,'' \textit{Journal of Spacecraft and Rockets}, vol. 56, no. 6, pp. 1662--1679, 2019.
  \item D. Y. Lee, H. Park, M. Romano, and J. Cutler, ``Development and experimental validation of a multi-algorithmic hybrid attitude determination and control system for a small satellite,'' \textit{Aerospace Science and Technology}, vol. 78, pp. 494--509, 2018.
  \item D. Y. Lee, R. Gupta, U. V. Kalabi\'{c}, S. Di Cairano, A. M. Bloch, J. W. Cutler, and I. V. Kolmanovsky, ``Geometric mechanics based nonlinear model predictive spacecraft attitude control with reaction wheels,'' \textit{Journal of Guidance, Control, and Dynamics}, vol. 40, no. 2, pp. 309--319, 2017.
  \item D. Y. Lee, J. W. Cutler, J. Mancewicz, and A. J. Ridley, ``Maximizing photovoltaic power generation of a space-dart configured satellite,'' \textit{Acta Astronautica}, vol. 111, pp. 283--299, 2015.
  \item J. T. Hwang, D. Y. Lee, J. W. Cutler, and J. R. R. A. Martins, ``Large-scale multidisciplinary optimization of a small satellite's design and operation,'' \textit{Journal of Spacecraft and Rockets}, vol. 51, no. 5, pp. 1648--1663, 2014.
  \item J. W. Choi, S. B. Lee, D. Y. Lee, U. S. Park, and Y. S. Suh, ``A fault isolation filter design using left eigenstructure assignment scheme,'' \textit{KSME International Journal}, vol. 14, no. 6, pp. 583--589, 2000.
\end{enumerate}

\subsubsection{Selected Conference Publications}

\begin{enumerate}
  \item Z. Luppen, M. Jacks, N. Baughman, M. Stilic, R. Nasers, B. Hertz, J. Cutler, D. Y. Lee, and K. Y. Rozier, ``Elucidation and analysis of specification patterns in aerospace system telemetry,'' in \textit{NASA Formal Methods Symposium}, pp. 527--537, 2022.
  \item Y. Lee and D. Y. Lee, ``Optimal trajectory generation for Mars atmospheric entry guidance using parameter optimization,'' in \textit{AIAA SCITECH 2022 Forum}, 2022.
  \item Z. A. Luppen, D. Y. Lee, and K. Y. Rozier, ``A case study in formal specification and runtime verification of a CubeSat communications system,'' in \textit{AIAA SciTech 2021 Forum}, Art. no. 0997, 2021.
  \item M. E. Nelson, D. Y. Lee, M. Kilcoin, L. Gordon, and W. Brown, ``Preparing CySat-1: A look at Iowa State University's first CubeSat,'' in \textit{SmallSat Conference}, 2020.
  \item M. Mammarella, D. Y. Lee, H. Park, E. Capello, M. Dentis, G. Guglieri, and M. Romano, ``Attitude control of a small spacecraft for Earth observation via tube-based robust model predictive control,'' in \textit{AIAA SPACE and Astronautics Forum and Exposition}, 2018.
  \item D. Y. Lee, S. Sharma, H. Park, and J. W. Cutler, ``Design and optimization of a small satellite communication system,'' in \textit{AIAA Aerospace Sciences Meeting}, 2018.
  \item D. Y. Lee, R. Gupta, U. V. Kalabi\'{c}, S. Di Cairano, A. M. Bloch, J. W. Cutler, and I. V. Kolmanovsky, ``Constrained attitude maneuvering of a spacecraft with reaction wheel assembly by nonlinear model predictive control,'' in \textit{2016 American Control Conference (ACC)}, pp. 4960--4965, 2016.
  \item J. Hwang, D. Y. Lee, J. Cutler, and J. R. R. A. Martins, ``Large-scale MDO of a small satellite using a novel framework for the solution of coupled systems and their derivatives,'' in \textit{54th AIAA/ASME/ASCE/AHS/ASC Structures, Structural Dynamics, and Materials Conference}, 2013.
  \item D. Lee, J. Springmann, S. Spangelo, and J. Cutler, ``Satellite dynamics simulator development using lie group variational integrator,'' in \textit{AIAA Modeling and Simulation Technologies Conference}, Art. no. 6430, 2011.
\end{enumerate}

\subsubsection{Synergistic Activities}

\begin{itemize}
  \item Design, development, and optimization of CubeSat fleet concepts.
  \item Spacecraft dynamics and control, including formation flight of CubeSat fleets.
  \item Mars entry, descent, and landing guidance.
  \item Nonlinear model predictive control and optimal control.
  \item Spacecraft systems and spacecraft dynamics instruction.
\end{itemize}
